\documentclass[12pt]{article}
\usepackage{amsmath, amssymb, amsthm}
\usepackage{geometry}
\geometry{margin=1in}

\newtheorem{exercise}{Exercise}
\theoremstyle{remark}
\newtheorem*{solution}{Solution}
\newtheorem*{remark}{Remark}

\title{Random Walks --- Week 2 Piazza Exercises}
\author{CUNY, Spring 2026}
\date{}

\begin{document}
\maketitle

\section*{Week 2: Properties of Simple Random Walk}

Throughout, $(S_n)_{n \ge 0}$ denotes the simple symmetric random walk (SRW) on $\mathbb{Z}$:
$S_0 = 0$ and $S_n = X_1 + \cdots + X_n$ where $X_1, X_2, \ldots$ are i.i.d.\ with
$P(X_i = +1) = P(X_i = -1) = \tfrac{1}{2}$.

% ----------------------------------------------------------------
\begin{exercise}[Permutation invariance vs.\ tail events]
Let $A_k = \{X_k = 1\}$. Determine whether $A_k$ is
(a) a tail event (i.e., $A_k \in \mathcal{T} = \bigcap_{n \ge 1} \sigma(X_{n+1}, X_{n+2}, \ldots)$), and
(b) permutation-invariant (i.e., $\pi^{-1}(A_k) = A_k$ for every finite permutation $\pi$ of $\mathbb{N}$).
\end{exercise}

\begin{solution}
\textbf{(a) $A_k$ is \emph{not} a tail event.} In fact $A_k \in \sigma(X_1, \ldots, X_k)$, so
$A_k$ is determined by finitely many coordinates and therefore cannot belong to the tail
$\sigma$-algebra $\mathcal{T} = \bigcap_{n} \sigma(X_{n+1}, X_{n+2}, \ldots)$.  To see this
more concretely, if $A_k \in \sigma(X_1,\ldots,X_j)$ for some $j < k$, then membership in
$A_k$ would be determined by the first $j$ coordinates alone, which is absurd: the element
$(1,0,\ldots)$ satisfies $\omega_k = 0$ if $k > 1$, giving a contradiction.

\textbf{(b) $A_k$ \emph{is} permutation-invariant.} For any finite permutation $\pi$, the
transformed sequence $(X_{\pi(1)}, X_{\pi(2)}, \ldots)$ has the same joint distribution as
$(X_1, X_2, \ldots)$ by the i.i.d.\ assumption. Hence $A_k = \{X_k = 1\}$ and
$\pi^{-1}(A_k) = \{X_{\pi^{-1}(k)} = 1\}$ agree up to permutation, and since the $X_i$
are i.i.d.\ the two events have equal probability. One can also write
\[
  A_k = \lim_{n \to \infty} \Bigl\{ \tfrac{1}{n}\sum_{j=1}^n X_j \in B \Bigr\}
\]
for a suitable Borel set $B$, exhibiting $A_k$ as a limit of symmetric functions, hence
permutation-invariant.
\end{solution}

% ----------------------------------------------------------------
\begin{exercise}[Stationary independent increments and Markov property]
Let $0 = k_0 < k_1 < k_2 < \cdots < k_n \le N$.
\begin{enumerate}
  \item[(1)] Show that $S_0,\; S_{k_1}-S_0,\; S_{k_2}-S_{k_1},\; \ldots,\; S_{k_n}-S_{k_{n-1}}$
             are independent.
  \item[(2)] Show that $S_n - S_k$ has the same distribution as $S_{n-k}$.
  \item[(3)] Show the Markov property:
             \[
               P(S_n = z_n \mid S_0=z_0, S_1=z_1, \ldots, S_{n-1}=z_{n-1})
               = P(S_n = z_n \mid S_{n-1} = z_{n-1}).
             \]
  \item[(4)] For $0 < k < n \le N$ and $P(S_k = z) > 0$,
             \[
               P(S_n = y \mid S_k = z) = P(S_{n-k} = y - z).
             \]
\end{enumerate}
\end{exercise}

\begin{solution}
Let $\alpha_i = k_i - k_{i-1}$ and $\beta_i = (\alpha_i + z_i)/2$ where
$z_i \in \mathbb{Z}$ denotes the value taken by the $i$-th increment.

\textbf{(1)} For the increment $S_{k_i} - S_{k_{i-1}} = X_{k_{i-1}+1} + \cdots + X_{k_i}$
to equal $z_i$, we need exactly $\beta_i$ of the $\alpha_i$ steps in that block to be $+1$.
This fixes $\beta_i$ of the coordinates in positions $\{k_{i-1}+1, \ldots, k_i\}$ in a
$\binom{\alpha_i}{\beta_i}$-many ways, leaving those in other blocks free.  Thus
\[
  P(S_{k_1}-S_0 = z_1, \ldots, S_{k_n}-S_{k_{n-1}}=z_n)
  = \prod_{i=1}^n \binom{\alpha_i}{\beta_i} 2^{-\alpha_i}
  = \prod_{i=1}^n P(S_{k_i} - S_{k_{i-1}} = z_i),
\]
establishing independence.

\textbf{(2)} Since the probability $P(X_{k+1}+\cdots+X_n = z) = \binom{n-k}{(n-k+z)/2} 2^{-(n-k)}$
depends only on the block length $n-k$ and not on which specific indices are used,
$S_n - S_k =^d S_{n-k}$.

\textbf{(3)} If $z_n \ne z_{n-1} \pm 1$, both sides are $0$. Otherwise,
\begin{align*}
  P(S_n=z_n \mid S_0=z_0,\ldots,S_{n-1}=z_{n-1})
  &= P(X_n = z_n - z_{n-1} \mid S_0=z_0,\ldots,S_{n-1}=z_{n-1}).
\end{align*}
Since $X_n$ is independent of $(S_0,\ldots,S_{n-1})$, the conditioning has no effect and
the probability equals $\tfrac{1}{2}$.  The same computation with conditioning only on
$S_{n-1}=z_{n-1}$ gives the same answer $\tfrac{1}{2}$, proving the Markov property.

\textbf{(4)} We compute directly:
\begin{align*}
  P(S_n=y \mid S_k=z)
  &= P(X_{k+1}+\cdots+X_n = y-z \mid X_1+\cdots+X_k = z) \\
  &= P(X_{k+1}+\cdots+X_n = y-z)  \quad \text{(independence of disjoint blocks)} \\
  &= P(S_{n-k} = y-z). \qquad \square
\end{align*}
\end{solution}

\begin{remark}
Property (1) says $S_n$ has \emph{independent increments}; Property (2) says it has
\emph{stationary increments}.  Together these two properties (plus right-continuity in the
continuous-time setting) characterize \emph{Lévy processes}.  Property (3) is the Markov
property; every process with independent increments is automatically Markov.
\end{remark}

% ----------------------------------------------------------------
\begin{exercise}[Return probability asymptotics]
\label{ex:return}
Show that
\[
  P(S_{2n}=0) - P(S_{2n+2}=0) = \frac{1}{2(n+1)}\,P(S_{2n}=0),
\]
and conclude $P(S_{2n}=0) - P(S_{2n+2}=0) \sim \dfrac{1}{2n\sqrt{\pi n}}$ as $n\to\infty$.
\end{exercise}

\begin{solution}
Using the exact formula $P(S_{2n}=0) = \binom{2n}{n}2^{-2n}$:
\begin{align*}
  P(S_{2n+2}=0) &= \binom{2n+2}{n+1}2^{-2n-2} \\
  &= \frac{(2n+2)!}{(n+1)!^2} \cdot \frac{1}{4^{n+1}} \\
  &= \frac{(2n+2)(2n+1)}{4(n+1)^2} \cdot \binom{2n}{n} 2^{-2n} \\
  &= \frac{2n+1}{2(n+1)} \cdot P(S_{2n}=0).
\end{align*}
Therefore
\[
  P(S_{2n}=0) - P(S_{2n+2}=0)
  = \Bigl(1 - \frac{2n+1}{2(n+1)}\Bigr)\,P(S_{2n}=0)
  = \frac{1}{2(n+1)}\,P(S_{2n}=0).
\]
By Stirling's formula, $P(S_{2n}=0) = \binom{2n}{n}2^{-2n} \sim \dfrac{1}{\sqrt{\pi n}}$, so
\[
  P(S_{2n}=0) - P(S_{2n+2}=0)
  \sim \frac{1}{2(n+1)} \cdot \frac{1}{\sqrt{\pi n}}
  \sim \frac{1}{2n\sqrt{\pi n}}
\]
since $n+1 \sim n$ as $n\to\infty$. \hfill$\square$
\end{solution}

% ----------------------------------------------------------------
\begin{exercise}[Reflection principle for SRW]
Let $a, b \in \mathbb{Z}$ with $a > 0$. Using the reflection principle, show that for all
$n \ge 1$ and $b \in \mathbb{Z}$,
\[
  P\!\bigl(S_n = b,\; \tau_a \le n\bigr) = P(S_n = 2a - b),
\]
where $\tau_a = \inf\{k \ge 0: S_k = a\}$.  Deduce the distribution of
$M_n = \max_{0 \le k \le n} S_k$:
\[
  P(M_n \ge a) = 2P(S_n \ge a) + P(S_n = a) \quad (a > 0,\; a \equiv n \pmod{2}).
\]
\end{exercise}

\begin{solution}
Consider any path $\omega = (\omega_0, \omega_1, \ldots, \omega_n)$ with
$\omega_n = b$ and $\tau_a \le n$ (i.e.\ the path reaches level $a$ by time $n$).
Define the \emph{reflected path} $\tilde\omega$ by leaving $\omega_k$ unchanged for
$k < \tau_a$ and setting $\tilde\omega_k = 2a - \omega_k$ for $k \ge \tau_a$.
Since each step is $\pm 1$, $\tilde\omega$ is a valid path ending at $2a - b$.

This reflection is a bijection between
$\{\text{paths reaching } a \text{ and ending at } b\}$ and
$\{\text{all paths ending at } 2a - b\}$,
so $P(S_n = b,\, \tau_a \le n) = P(S_n = 2a - b)$.

For the maximum: $\{M_n \ge a\}$ contains paths that reach $a$ by time $n$.  Conditioning
on the value of $S_n$ and splitting into $b \ge a$ (path still at or above $a$) and
$b < a$ (path reflected back) gives
\[
  P(M_n \ge a) = P(S_n \ge a) + P(S_n = 2a - S_n \le a - 1)
  = 2P(S_n \ge a) + P(S_n = a),
\]
where the $P(S_n = a)$ term accounts for paths that hit $a$ exactly at time $n$.
\hfill$\square$
\end{solution}

% ----------------------------------------------------------------
\begin{exercise}[Distribution of the first passage time $\tau_a$]
For $a \ge 1$, show that
\[
  P(\tau_a = n) = \frac{a}{n}\,P(S_n = a), \qquad n \ge 1.
\]
\end{exercise}

\begin{solution}
We use the identity (derivable from the ballot problem or directly from the reflection
principle):
\[
  P(\tau_a = n)
  = \tfrac{1}{2}P(S_{n-1} = a-1) - \tfrac{1}{2}P(S_{n-1} = a+1).
\]
Set $k = (a+n)/2$ (assuming $n \equiv a \pmod 2$, otherwise both sides are $0$). Then
\begin{align*}
  P(\tau_a = n)
  &= \tfrac{1}{2} \cdot 2^{-(n-1)} \Bigl[\binom{n-1}{k-1} - \binom{n-1}{k}\Bigr] \\
  &= 2^{-n} \Bigl[\frac{k}{n}\binom{n}{k} - \frac{n-k}{n}\binom{n}{k}\Bigr] \\
  &= 2^{-n} \cdot \frac{2k-n}{n} \cdot \binom{n}{k}.
\end{align*}
Since $2k - n = 2\cdot\frac{a+n}{2} - n = a$, we obtain
\[
  P(\tau_a = n) = \frac{a}{n} \cdot \binom{n}{(a+n)/2} \cdot 2^{-n}
  = \frac{a}{n}\,P(S_n = a). \qquad\square
\]
This identity is a special case of the \emph{Hitting-Time Theorem} (ballot theorem for
integer-valued walks with increments $\le 1$).
\end{solution}

\end{document}
